\cotefile{77}
\scanpage{66}
% archive sheet number: 65
\uline{Discriminant de formes quadratiques.}

Soit $I$ ens d'indices fini, $\mathscr A_I=\mathbb Z[I\times I]$ l'anneau \inserted{universelle} \deleted{\unsure{univ.}} \inserted{\unsure{le|les|des}} matrices carrées de type $I$, $\mathscr A\unsure{Q|S}_I=\mathbb Z[\operatorname{Sym}^2(I)]$ \deleted{le groupe des} \inserted{celui de la} matrices carrées symétriques universelle, \deleted{de} dont les générateurs sur $Z$ sont $\deleted{{\unsure{\mathscr A|\mathbb Z|A\mathbb Z}}^{I\times I}\longrightarrow Z}$ de la forme $\deleted{X}Y_J$, où $J$ parcourt l'ens des parties de $I$ de cardinal $1$ ou $2$. \unsure{Écriv.|Écr.|Coord.} ds $\mathscr A_I$ $\deleted{\mathscr A Q_I\longleftarrow\mathscr A Q_J}$ les $X_{ij}$\inserted{$(i,j\in I\times I)$}. \deleted{\unsure{les}} considérons
\[
\mathscr A_I\xrightarrow{\varphi}\mathscr A Q_I
\]
défini par
\[
\varphi(X_{\unsure{\{i,j|i,j|ij}})=
\begin{cases}
\deleted{X}Y_{\{i,j\}}&\text{si }i\ne j\\
2\deleted{X}Y_{\{i,j\}}=2X_{\{i\}}&\text{si }i=j.
\end{cases}
\]
qui \deleted{associ\unsure{e}} correspond à l'hom de foncteurs
\[
MS_I(k)\longrightarrow M_I(k)
\]
associant à une forme quadratique de coeff. les $a_{\{i,j\}}$, la matrice de celle-ci.

Soit $\Delta_0\in\mathscr A_I$ le déterminant de la $I$-matrice universelle, $\Delta=\varphi(\Delta_0)$ son image, qui est le discriminant de la forme quadratique universelle indexée par $I$.

On a
\[
\Delta_0=\sum_{\sigma\in\unsure{\mathfrak S|\mathscr S}_I}\varepsilon_\sigma
\overbrace{\prod_{i\in I}X_{i,\sigma i}}^{X_\sigma}
\]
Considérons la décomposition de $I$ en cycles suivant $\sigma$, le produit s'écrit -- en mettant à
